% =====================================================================
%  Resumo do projeto de dissertação - Seminário de Dissertação I
%  Limite: 800 palavras. Compilar no Overleaf com pdfLaTeX.
% =====================================================================
\documentclass[12pt,a4paper]{article}
\usepackage[utf8]{inputenc}
\usepackage[T1]{fontenc}
\usepackage[brazil]{babel}
\usepackage{mathptmx}
\usepackage[top=2.5cm, left=3cm, bottom=2cm, right=2cm]{geometry}
\usepackage{setspace}
\setstretch{1.15}
\usepackage{indentfirst}
\setlength{\parindent}{1.25cm}
\usepackage{microtype}
\usepackage{titlesec}
\titleformat{\section}{\normalfont\bfseries}{}{0em}{}
\titlespacing*{\section}{0pt}{8pt}{3pt}
\pagestyle{empty}

\begin{document}

\begin{center}
  {\bfseries O EFEITO DA COP30 SOBRE O EMPREGO FORMAL EM ALOJAMENTO E ALIMENTAÇÃO EM BELÉM}\\[4pt]
  Marcela Paz da Silva Enriquez\\
  Resumo do projeto de dissertação -- Seminário de Dissertação I -- PPGE/UFF
\end{center}

\section{Contexto e pergunta}

Grandes eventos costumam ser defendidos pela promessa de gerar empregos. Foi assim com a COP30,
realizada em Belém em novembro de 2025, a primeira COP sediada na Amazônia.
A literatura internacional sobre megaeventos, porém, encontra efeitos sobre o emprego pequenos,
concentrados em poucos setores e passageiros. A
pesquisa pergunta: \textbf{a COP30 mudou o emprego formal em Alojamento e Alimentação em Belém, e esse
efeito permaneceu depois que o evento terminou?} Como hotéis, restaurantes e bares são os primeiros a
receber o gasto dos visitantes, a análise do emprego se concentra nesses setores.

\section{Objetivos}

O objetivo geral é estimar o efeito causal da COP30 sobre o saldo de empregos formais do setor em
Belém. Os objetivos específicos são: (i) separar os efeitos na preparação, no mês do evento e na
desmobilização; (ii) verificar se o efeito se concentra no alojamento ou na alimentação; e (iii)
verificar se o emprego se espalhou para a Região Metropolitana.

\section{Hipóteses}

H1: houve contratação antecipada na fase de preparação, depois do anúncio do evento. H2: o saldo
aumentou no mês do evento, puxado pelo alojamento. H3: na desmobilização, o saldo ficou abaixo do
que seria esperado sem o evento. H4: o efeito foi temporário.

\section{Dados}

A base usa os microdados do Novo CAGED, que registram cada contratação e demissão com carteira
assinada. Para cada capital e cada mês, de janeiro de 2022 a julho de 2026, calcula-se o saldo (contratações
menos demissões) nas divisões 55 e 56 da CNAE. O saldo é dividido pela média mensal de contratações da própria cidade antes do evento. O resultado é
um painel balanceado de 26 capitais observadas por 55 meses. O Rio de Janeiro é excluído porque sediou
a Cúpula do G20 em novembro de 2024, um evento do mesmo tipo.

\section{Metodologia}

O problema central é que não se observa o emprego que Belém teria tido sem a COP30, o contrafactual.
Comparar Belém antes e depois não basta, porque o emprego também muda com a estação e a conjuntura. A estratégia é a de \textbf{diferenças em diferenças}: compara-se a
variação do emprego em Belém com a variação nas demais capitais no mesmo período. A variação das outras
capitais representa o que teria acontecido com Belém sem o evento, e a diferença entre as duas
variações é o efeito da COP30. O modelo é estimado por mínimos quadrados ordinários com efeitos fixos
de capital e de mês. O desenho segue Leoni e Boto-García (2023), que avaliam um
choque temporário sobre a hotelaria com uma única unidade tratada.

A hipótese de identificação é a de tendências paralelas: sem a COP30, Belém teria seguido a mesma
tendência das outras capitais. Como há uma única
cidade tratada, a significância é avaliada por placebo, tratando cada capital como se tivesse sediado
a COP30. Os resultados são
conferidos pelo controle sintético (Abadie, Diamond e Hainmueller, 2010) e por testes que isolam o
Círio de Nazaré e a conferência realizada em Belém em novembro de 2024.
Uma análise complementar aplica o mesmo desenho aos municípios do Pará, comparando Belém e sua Região
Metropolitana com o interior, para medir se o emprego transbordou para os vizinhos.

\section{Resultados esperados e contribuição}

Espera-se um efeito positivo e concentrado no alojamento durante o evento, seguido de desaceleração
na desmobilização. Essa expectativa se apoia em estudos que avaliaram os efeitos econômicos de
outros megaeventos: a Copa do Mundo de 2006 na Alemanha (Hagn e Maennig, 2009), a Copa do Mundo de
2014 no Brasil (Baumann e Matheson, 2018), os Jogos Olímpicos de Atlanta em 1996 (Feddersen e
Maennig, 2013) e o Jubileu de 2000 em Roma (Bronzini, Mocetti e Mongardini, 2019). Em geral, esses
estudos encontram efeitos modestos, concentrados no turismo e no período do evento. Dados preliminares
do Pará mostram saldo positivo depois da COP30, mas menor que nos anos anteriores. A pesquisa contribui com a primeira avaliação \textit{ex post} de uma conferência do clima sobre o
mercado de trabalho, com dados administrativos brasileiros. Uma limitação é que o CAGED capta apenas o emprego formal,
relevante em uma região de alta informalidade.

\section{Referências}
\small
ABADIE, A.; DIAMOND, A.; HAINMUELLER, J. Synthetic control methods for comparative case studies.
\textit{Journal of the American Statistical Association}, v. 105, n. 490, p. 493--505, 2010.\par
BAUMANN, R.; MATHESON, V. Mega-events and tourism: the case of Brazil. \textit{Contemporary Economic
Policy}, v. 36, n. 2, p. 292--301, 2018.\par
BRONZINI, R.; MOCETTI, S.; MONGARDINI, M. \textit{The economic effects of big events}: evidence from the
Great Jubilee 2000 in Rome. Roma: Banca d'Italia, 2019. (Temi di discussione, n. 1208).\par
FEDDERSEN, A.; MAENNIG, W. Mega-events and sectoral employment: the case of the 1996 Olympic Games.
\textit{Contemporary Economic Policy}, v. 31, n. 3, p. 580--603, 2013.\par
HAGN, F.; MAENNIG, W. Large sport events and unemployment: the case of the 2006 soccer World Cup in
Germany. \textit{Applied Economics}, v. 41, n. 25, p. 3295--3302, 2009.\par
LEONI, V.; BOTO-GARCÍA, D. The effect of natural disasters on hotel demand, supply and labour markets:
evidence from the La Palma volcano eruption. \textit{Environmental and Resource Economics}, v. 86,
n. 4, p. 755--780, 2023.\par

\end{document}
